% Part: first-order-logic
% Chapter: model-theory
% Section: overspill

\documentclass[../../../include/open-logic-section]{subfiles}

\begin{document}

\olfileid[hi]{mod}{bas}{ove}
\olsection{अतिप्रवाह}

\begin{thm}
\ollabel{overspill} यदि वाक्यों के किसी समुच्चय $\Gamma$ के मनमाने रूप से
बड़े परिमित प्रतिरूप हैं, तो उसका एक अनंत प्रतिरूप भी है।
\end{thm}

\begin{proof}
$\Gamma$ की भाषा में गणनीय रूप से अनेक नए नियतांक $c_0$, $c_1$, \dots
जोड़कर उसे विस्तारित करें और समुच्चय $\Gamma \cup \{c_i \neq c_j :
i \neq j\}$ पर विचार करें। यह कहना कि $\Gamma$ के मनमाने रूप से बड़े परिमित
प्रतिरूप हैं, इसका अर्थ है कि प्रत्येक $m >0$ के लिए कोई $n\ge m$ ऐसा है कि
$\Gamma$ के किसी प्रतिरूप के प्रांत में~$n$ अवयव हैं। इससे $\Gamma \cup \{c_i \neq
c_j : i \neq j\}$ परिमिततः संतुष्टनीय है। संहतता के अनुसार, $\Gamma
\cup \{c_i \neq c_j : i \neq j\}$ का कोई प्रतिरूप $\Struct M$ है जिसका प्रांत
अनंत होना ही चाहिए, क्योंकि वह सभी असमानताओं $c_i \neq c_j$ को संतुष्ट करता
है।
\end{proof}

\begin{prop}
\ollabel{inf-not-fo} 
किसी भी प्रथम-कोटि भाषा का कोई ऐसा वाक्य $!A$ नहीं है जो किसी
!!a{structure}~$\Struct M$ में तभी और केवल तभी सत्य हो जब उस !!{structure}
का प्रांत $\Domain{M}$ अनंत हो।
\end{prop}

\begin{proof}
यदि ऐसा कोई $!A$ होता, तो उसका निषेध $\lnot !A$ सभी और केवल परिमित
!!{structure}s में सत्य होता। अतः उसके मनमाने रूप से बड़े परिमित प्रतिरूप
होते, किंतु कोई अनंत प्रतिरूप न होता; यह \olref{overspill} के विरुद्ध है।
\end{proof}

\end{document}
